% Pascal Michaillat, Problem Set on Dynamic Programming. CC BY 4.0.
% Adapted for tp-latex; see cover.tex and the editorial note below.
\documentclass{article}
\usepackage{hyperref}
\begin{document}
\section*{Problem Set on Dynamic Programming}
Pascal Michaillat. \href{https://github.com/pmichaillat/math-for-macro/blob/main/ps1.tex}{Original LaTeX source}.

\subsection*{Problem 1}

Consider the following optimal growth problem: Given initial capital $k_{0}>0$, choose consumption $\left\lbrace c_{t}\right\rbrace  _{t =0}^{\infty}$ to maximize utility
\[

\sum_{t=0}^{\infty}\beta^{t} \ln\left(c_{t}\right)

\]
subject to the resource constraint
\[

k_{t+1}=a k_{t}^{\alpha}-c_{t}.

\]
The parameters satisfy $0<\beta<1$, $a>0$, $0<\alpha <1$.

\wcOl[type=A]{
\wcLi{Derive the optimal law of motion of consumption $c_{t}$ using a Lagrangian.}
\wcLi{Identify the state variable and the control variable.}
\wcLi{Write down the Bellman equation.}
\wcLi{Derive the following Euler equation: 
\[

c_{t+1}=\beta  \alpha  a k_{t+1}^{\alpha -1} c_{t}.

\]}
\wcLi{Derive the first two value functions, $V_{1}(k)$ and  $V_{2}(k)$, obtained by iteration on the Bellman equation starting with the value function $V_{0}(k) \equiv 0$.}
\wcLi{The process of determining the value function by iterations using the Bellman equation is commonly used to solve dynamic programs numerically. The algorithm is called \textit{value function iteration}. For this optimal growth problem, one can show using value function iteration that the value function is
\[V(k) =\kappa +\frac{\ln\left(k^{\alpha}\right)}{1-\alpha \beta},\]
where $\kappa$ is a constant. Using the Bellman equation, determine the policy function $k'(k)$ associated with this value function.}
\wcLi{In light of these results, for which reasons would you prefer to use the dynamic-programming approach instead of the Lagrangian approach to solve the optimal growth problem? And for which reasons would you prefer to use the Lagrangian approach instead of the dynamic-programming approach?}
}

\subsection*{Problem 2}

Consider the problem of choosing consumption $\left\lbrace c_{t}\right\rbrace _{t=0}^{\infty}$ to maximize expected utility
\[

\mathbb{E}_{0}\left(\sum_{t=0}^{\infty}\beta^{t} u(c_{t})\right)

\]
subject to the budget constraint
\[

c_{t}+p_{t} s_{t+1}=(d_{t}+p_{t}) s_{t}.

\]
$d_{t}$ is the dividend paid out for one share of the asset, $p_{t}$ is the price of one share of the asset, and $s_{t}$ is the number of shares of the asset held at the beginning of period $t$. In equilibrium, the price $p_{t}$ of one share is solely a function of dividends $d_{t}$. Dividends can only take two values $d_{l}$ and $d_{h}$, with $0<d_{l}<d_{h}$. Dividends follow a Markov process with transition probabilities 
\[

\mathbb{P}\left(d_{t+1}=d_{l}\mid d_{t}=d_{l}\right) =\mathbb{P}\left(d_{t+1}=d_{h}\mid d_{t}=d_{h}\right) =\rho

\]
with $0.5<\rho 1$.

\wcOl[type=A]{
\wcLi{Identify state and control variables.}
\wcLi{Write down the Bellman equation.}
\wcLi{Derive the following Euler equation: 
\[

p_{t} u'(c_{t}) =\beta  \mathbb{E}\left((d_{t+1}+p_{t+1})  u'(c_{t+1}) \mid d_{t}\right) .

\]}
\wcLi{Suppose that $u(c) =c$. Show that the asset price is higher when the current dividend is high.}
}

\subsection*{Problem 3}

Consider the following optimal growth problem: Given initial capital $k_{0}>0$, choose consumption and labor $\left\lbrace c_{t},l_{t}\right\rbrace _{t=0}^{\infty}$ to maximize utility
\[

\sum_{t=0}^{\infty}\beta^{t} u(c_{t},l_{t})

\]
subject to the law of motion of capital
\[
\begin{aligned}
k_{t+1}&=a_{t} f(k_{t},l_{t}) -c_{t}.
\end{aligned}
\]
In addition, we impose $0\le l_{t}\le 1$. The discount factor $\beta \in (0,1)$. The function $f$ is increasing and concave in both arguments. The function $u$ is increasing and concave in $c$, decreasing and convex in $l$. 

\paragraph{Deterministic case} First, suppose $a_{t}=1$ for all $t$.

\wcOl[type=A]{
\wcLi{What are the state and control variables?}
\wcLi{Write down the Bellman equation.}
\wcLi{Derive the following optimality conditions: 
\[
\begin{aligned}
\frac{\partial u(c_{t},l_{t})}{\partial c_{t}} &=\beta  \frac{\partial u(c_{t+1},l_{t+1})}{\partial c_{t+1}} \cdot \frac{\partial f(k_{t+1},l_{t+1})}{\partial k_{t+1}}\\
\frac{\partial u(c_{t},l_{t})}{\partial c_{t}} \cdot \frac{\partial f(k_{t},l_{t})}{\partial l_{t}} &=-\frac{\partial u(c_{t},l_{t})}{\partial l_{t}}.
\end{aligned}
\]}
\wcLi{Suppose that the production function $f\left(k,l\right) =k^{\alpha}\cdot l^{1-\alpha}$. Determine the ratios $c/k$ and $l/k$ in steady state.}
}

\paragraph{Stochastic case} Now, suppose $a_{t}$ is a stochastic process that takes values $a_{1}$ and $a_{2}$ with the following transition probabilities: $\mathbb{P}\left(a_{t+1}=a_{1}\mid a_{t}=a_{1}\right) =\mathbb{P}\left(a_{t+1}=a_{2}\mid a_{t}=a_{2}\right) =\rho.$
\wcOl[type=A,start=5]{
\wcLi{Write down the Bellman equation.}
\wcLi{Derive the optimality conditions.}
}


\subsection*{Editorial note}
The source writes the bound on the transition probability in Problem 2 as $0.5<\rho 1$. This apparent typographical error is retained here; the expected upper-bound condition is $\rho<1$. The original question labels and their continuation from D to E in Problem 3 are preserved. Suggested solutions prepared for this adaptation are available on a separate page. They are not an official answer key by the original author.
\href{solutions.tex}{Read the suggested solutions}
\end{document}
